\section{Understanding Large Language Models}

\subsection{Oque caracteriza uma LLM?}
Um Large Language Model, também conhecido como LLM, é um modelo de inteligência criado a partir de algoritmos matemáticos, capaz de realizar a tarefa para a qual foi designado. É um modelo de inteligência artificial baseado em aprendizado de máquina, no qual recebe dados, como textos, artigos, notícias, entre outras informações disponíveis na internet, que o tornam capaz de trabalhar com a linguagem humana. Esses "sistemas inteligentes" não são capazes de compreender o nosso idioma da mesma forma que nós, mas são capazes de se comunicar através dele.

O que categoriza uma inteligência artificial (IA) como um LLM começa pelo próprio conceito da sigla. O "Large" vem da grande escala de informações disponibilizadas para que um modelo de inteligência artificial "aprenda" a linguagem humana. A princípio, isso envolve a capacidade da inteligência artificial de prever a próxima palavra com base nas palavras anteriores.

\subsection{Qual a Diferença entre IA, Machine Learning e Deep Learning?}
O princípio de uma inteligência artificial é a capacidade de um sistema realizar tarefas que necessitariam de capacidades inteligentes (como pensar ou até mesmo escrever), utilizando algoritmos capazes de detectar padrões a partir de dados.

Já o Machine Learning, subcampo da IA, é um modelo de inteligência artificial que não depende somente de algoritmos pré-programados, mas é também capaz de "aprender" a partir de dados para realizar uma tarefa em específico. É necessária a intervenção humana para a caracterização dos dados.

O Deep Learning é um subcampo do Machine Learning, onde o algoritmo do modelo possui mais camadas de processamento de informações e dados (redes neurais), tornando suas respostas mais coerentes com o assunto abordado e, dependendo da situação em que for aplicado, mais precisas.

\subsection{Como os modelos GPT estão organizados?}
Os modelos GPT são baseados na arquitetura Transformer e utilizam uma estrutura do tipo decoder-only. Eles processam o texto como uma sequência de tokens e, a partir do contexto anterior, aprendem a prever o próximo token.

\subsection{Quais são as principais etapas do treinamento de um LLM?}
A LLM possui 2 principais etapas de treinamento que são, pre training e o fine Tunning, durante a fase de pre training o modelo recebe uma vasta quantidades de dados possibilitando ele ser capaz de "aprender" a linguagem humana, sendo capaz de detectar padrões e prever a próxima palavra / token em sequencia. 

Após o pré treinamento é realizado a etapa do Fine Tunning, aonde o modelo recebe dados numa proporção menor, porém com um determinado foco, direcionando e treinando o modelo para ter um melhor desempenho para uma tarefa especifica. nesse Fine Tunning consiste 2 principais tipos:
\begin{itemize}
    \item Classification Fine Tunning: que consiste em fazer o modelo ser capaz de classificar textos em diversas categorias específicas, como por exemplo no livro é tratado os e-mails de spam e não spam.
    \item Instruction Fine Tunning: consiste em fazer o modelo ser capaz de "compreender" instruções e conseguir gerar respostas adequadas com oque foi solicitado.
\end{itemize}

\subsection{Qual é o fluxo geral de funcionamento de um modelo de linguagem?}
O modelo recebe um texto, transforma-o em tokens, processa esses tokens para identificar padrões e prevê o próximo token. Esse processo é repetido várias vezes até formar o texto de saída.

Exemplo:
Texto → Tokenização → Tokens → Modelo de linguagem → Previsão do próximo token → Repetição → Texto gerado
